\section[Memulai: Instalasi PHP dan Environment]{Memulai: Instalasi PHP dan Environment (XAMPP, Server Lokal)}

Untuk menjalankan kode PHP, Anda memerlukan sebuah lingkungan server karena PHP adalah bahasa \textit{server-side} yang membutuhkan web server dan interpreter PHP. Cara termudah bagi pemula adalah menggunakan paket \textbf{XAMPP}, yang menyediakan Apache (web server), MySQL/MariaDB (basis data), PHP, dan Perl dalam satu instalasi lengkap. XAMPP tersedia untuk Windows, Linux, dan macOS, sehingga sangat fleksibel digunakan di berbagai sistem operasi \cite{phpManual}.

Proses instalasi XAMPP sangat sederhana: unduh installer dari situs resmi Apache Friends, jalankan installer, lalu pilih komponen yang ingin dipasang (minimal Apache dan PHP). Setelah instalasi selesai, buka XAMPP Control Panel dan klik \textbf{Start} pada modul Apache. Jika berhasil, modul Apache akan berwarna hijau, menandakan web server sudah aktif. Semua file PHP yang ingin Anda jalankan harus disimpan di dalam folder \code{htdocs} (biasanya di \code{C:\textbackslash xampp\textbackslash htdocs} pada Windows), lalu diakses melalui browser dengan URL \code{http://localhost/nama-file.php} \cite{w3schoolsPhp}.

Langkah pertama setelah instalasi adalah memverifikasi bahwa PHP berjalan dengan benar menggunakan fungsi \code{phpinfo()}. Buat file baru bernama \code{info.php} di folder \code{htdocs}, lalu isi dengan kode berikut. Halaman yang dihasilkan akan menampilkan seluruh informasi konfigurasi PHP: versi PHP, ekstensi yang aktif, pengaturan server, dan banyak lagi. Ini sangat berguna untuk memeriksa apakah lingkungan pengembangan sudah siap.

\begin{webcode}[language=PHP, caption={Mengecek instalasi PHP dengan phpinfo()}]
<?php
// File: htdocs/info.php
// Akses via: http://localhost/info.php
phpinfo();
?>
\end{webcode}

\begin{konsep}
\textbf{Struktur Folder htdocs:} Setiap folder di dalam \code{htdocs} menjadi path di URL. Misalnya, file \code{htdocs/proyek/index.php} diakses dengan \code{http://localhost/proyek/index.php}. Memahami pemetaan folder ke URL ini penting untuk mengorganisasi proyek web Anda.
\end{konsep}

Selain XAMPP, terdapat beberapa alternatif lingkungan pengembangan PHP yang dapat digunakan sesuai preferensi dan sistem operasi. Masing-masing memiliki kelebihan dan kekurangan, namun semuanya menyediakan fungsi dasar yang sama: menjalankan web server dan interpreter PHP.

\begin{table}[ht]
\centering
\begin{tabular}{|P{2.5cm}|P{2.5cm}|P{4cm}|P{3cm}|}
\hline
\textbf{Paket} & \textbf{Platform} & \textbf{Komponen} & \textbf{Keterangan} \\
\hline
XAMPP & Windows, Linux, macOS & Apache, MySQL, PHP, Perl & Paling populer, lengkap \\
\hline
WAMP & Windows & Apache, MySQL, PHP & Khusus Windows \\
\hline
LAMP & Linux & Apache, MySQL, PHP & Instalasi manual via terminal \\
\hline
MAMP & macOS, Windows & Apache, MySQL, PHP & Populer untuk macOS \\
\hline
PHP Built-in Server & Semua OS & PHP saja & Ringan, tanpa Apache \\
\hline
\end{tabular}
\caption{Perbandingan paket server lokal untuk pengembangan PHP}
\end{table}

PHP juga menyediakan web server bawaan yang sangat praktis untuk pengembangan dan pengujian cepat tanpa perlu menginstal Apache. Cukup buka terminal, navigasi ke folder proyek, lalu jalankan perintah \code{php -S localhost:8000}. Server bawaan ini akan melayani file di folder saat ini dan menampilkan log request di terminal. Namun, server ini \textbf{tidak} disarankan untuk lingkungan produksi karena hanya dapat menangani satu request secara bersamaan dan tidak memiliki fitur keamanan lengkap \cite{phpRightWay}.

\begin{webcode}[language=Bash, caption={Menjalankan PHP built-in server}]
# Navigasi ke folder proyek
cd /path/to/proyek

# Jalankan server di port 8000
php -S localhost:8000

# Akses via browser: http://localhost:8000
\end{webcode}

